\section{Условие}

Цель лабораторной работы — познакомиться с построением простых трёхмерных объектов, их отображением на двумерной плоскости и интерактивным изменением проекции и ориентации фигуры. Работа выполнена на основе предоставленного проекта Vulkan Starter App на языке C++20.

Согласно варианту № 4 необходимо построить усечённый правильный тетраэдр, гранями которого являются правильные треугольники и правильные шестиугольники. Особое внимание уделяется правильности геометрии: грани должны быть плоскими, а длины всех рёбер — одинаковыми.

Программа должна использовать Vulkan, GLFW и Dear ImGui, работать в реальном времени и позволять пользователю динамически менять проекцию и трансформации объекта. Отображение фигуры должно учитывать выбранную проекцию и обеспечивать взаимодействие через интерфейс и устройства ввода.

В реализованном приложении предусмотрены перспективная и ортографическая проекции, поворот вокруг трёх координатных осей, изменение масштаба отображения, каркасный режим и автоматическое вращение. Переключение режимов выполняется без перезапуска программы.

В задании отдельно перечислены дополнительные возможности для повышения оценки. В данном отчёте рассматривается базовый вариант и реализованное переключение проекций. Перенос и независимое растяжение объекта по осям, сложная траектория движения, настройка цвета и процедурная окраска вершин, а также несколько объектов с отдельными наборами дескрипторов не реализованы. Автоматическое вращение не является выполнением задания о сложной траектории.
